\documentclass[11pt,letterpaper]{article}
\usepackage[T1]{fontenc}
\usepackage[utf8]{inputenc}
\usepackage{lmodern,microtype}
\usepackage[margin=1in]{geometry}
\usepackage{amsmath,amssymb,amsthm,mathtools}
\usepackage{booktabs,longtable,array}
\usepackage{enumitem,listings,xcolor}
\usepackage[hidelinks]{hyperref}
\hypersetup{pdftitle={Speeding Up Exact Unknot Recognition},
  pdfauthor={fastunknot 0.2.0 technical report},
  pdfsubject={Exact modular filters, sparse Khovanov cancellation, and connected-sum factors},
  pdfkeywords={unknot recognition, Khovanov homology, exact computation, complexity}}
\usepackage{fancyhdr}
\pagestyle{fancy}\fancyhf{}
\fancyhead[L]{\small Exact unknot recognition: implementation and analysis}
\fancyhead[R]{\small fastunknot 0.2.0}
\fancyfoot[C]{\thepage}
\setlength{\headheight}{14pt}
\setlength{\emergencystretch}{2em}
\setlist{itemsep=2pt,topsep=4pt}
\lstset{basicstyle=\small\ttfamily,columns=fullflexible,keepspaces=true,
  breaklines=true,frame=single,framerule=0.2pt,showstringspaces=false,
  xleftmargin=3pt,xrightmargin=3pt,aboveskip=9pt,belowskip=9pt}
\newtheorem{theorem}{Theorem}[section]
\newtheorem{proposition}[theorem]{Proposition}
\newtheorem{lemma}[theorem]{Lemma}
\newtheorem{corollary}[theorem]{Corollary}
\theoremstyle{definition}\newtheorem{definition}[theorem]{Definition}
\theoremstyle{remark}\newtheorem{remark}[theorem]{Remark}
\newcommand{\F}{\mathbb F}
\newcommand{\Z}{\mathbb Z}
\newcommand{\Q}{\mathbb Q}
\newcommand{\Kh}{\operatorname{Kh}}
\newcommand{\rKh}{\widetilde{\operatorname{Kh}}}
\newcommand{\poly}{\operatorname{poly}}
\newcommand{\code}[1]{\texttt{\detokenize{#1}}}
\newcommand{\UNKNOT}{\textsc{Unknot}}
\newcommand{\KNOTTED}{\textsc{Knotted}}
\newcommand{\UNKNOWN}{\textsc{Unknown}}
\title{Speeding Up Exact Unknot Recognition\\[4pt]
\large One-sided modular filters, sparse cancellation, and compact connected-sum factors}
\author{Implementation and technical analysis accompanying fastunknot 0.2.0}
\date{18 September 2026}
\begin{document}
\maketitle
\begin{abstract}
We improve the Python implementation supplied in the \code{fast/} directory of
\code{Knots.zip}, preserving an exact and complete Khovanov-based fallback.
The main changes are exact finite-field Alexander and Jones rejection filters,
immutable memoized cobordism arithmetic, lazy fill-aware cancellation, faster
scan ordering and descending tests, and certified diagrammatic connected-sum
factorization. On the supplied 36-crossing random five-braid, the original
engine exceeds a 60-second same-machine budget; the optimized unfactored rank
computation completes in 1.110 seconds. With caching but the old LIFO pivot
strategy, the same task takes 49.576 seconds. Default recognition of the Conway
and Kinoshita--Terasaka examples is respectively 21.1 and 24.5 times faster.

The general worst-case bound remains $2^{O(n)}$. We do not claim to implement
Lackenby's general quasi-polynomial algorithm. A proved improvement is
$\poly(n)+\sum_i 2^{O(n_i)}$, where $n_i$ are the sizes of the diagrammatic
connected-sum factors actually found. Consequently the implementation is
polynomial when their maximum size is $O(\log n)$ and quasi-polynomial when it is
$O(\log^2 n)$. A sum of 100 trefoils is processed in 65.6 milliseconds, with at
most 18 live objects, while its reduced rank $3^{100}$ is represented exactly.
The deliverable includes the original baseline, runnable source, certificates,
30 regression-test methods, and an exhaustive comparison on 2856 small closed
braid words against structurally independent computations, with no mismatches.
\end{abstract}
\newpage
\setcounter{tocdepth}{2}
{\small\tableofcontents}
\newpage

\section{Outcome, scope, and the supplied baseline}

The objective was first to pursue the quasi-polynomial target
$n^{O(\log n)}$, and otherwise to improve the supplied recognizer in a verifiable
way. The delivered program achieves the latter, together with a stronger bound
on an explicitly characterized class of diagrams. There is no unimplemented
oracle on its execution path. It uses only Python's standard library and runs
without network access or an external topology package.

Throughout, $n$ is the number of crossings in a validated, one-component,
classical planar diagram. Edge labels are compacted; after compaction the input
has $O(n\log n)$ bits. Reading arbitrarily long user-supplied integer labels is
an additional input-encoding cost, not a function of $n$ alone. Complexity bounds
assume sufficient memory and no user-imposed timeout. The resource-limited
interface deliberately permits the third answer \UNKNOWN.

The original program is preserved under \code{baseline/}; the improved version
is under \code{fast/}. The original program already had several important
components: planar-diagram, braid, and grid input; checked decreasing R1/R2
moves; a descending-diagram test; an exact Alexander polynomial; and a scanning
Khovanov calculation over $\F_2$ using local dotted cobordisms. This is a useful
starting point, not a naive full-cube enumerator. The improvement therefore
concentrates on avoiding unnecessary homology computations, reducing repeated
algebra, reducing cancellation fill-in, and avoiding explicit tensor-product
multiplicities.

The supplied historical benchmark was produced on a different Windows/Python
installation. It records a 600-second timeout on the large five-braid. That
number is useful for identifying a difficult test, but is \emph{not} used in any
same-machine speedup ratio in this report. The new measurements and their
individual samples are in \code{results/benchmark.json}. Input and source hashes
are provided so that the comparison can be repeated.

\subsection{What is and is not changed}

The topological meaning of accepted inputs is unchanged. The recognizer still
uses exact arithmetic. The final decision criterion is still total reduced
$\F_2$ Khovanov rank one. It does not compute integral torsion or the full quantum
bigrading. The returned degree histogram is in \emph{unshifted cube degree};
this distinction matters when comparing diagrams related by Reidemeister moves.

The default execution method may now report a modular Alexander or Jones
rejection instead of a full polynomial or homology computation. These are exact
one-sided implications, not randomized guesses. The low-level rank API now also
validates the PD and any explicitly supplied crossing order. The writhe
implementation is corrected to use both oriented branches, while the original
local signs used in the Fox matrix retain their original convention.

\section{What the Lackenby material does and does not supply}

The supplied talk announces a general quasi-polynomial unknot-recognition
algorithm; its theorem slide attributes the result to 2021. Its later slides
identify four acceleration ingredients: surfaces of quadratic pattern
complexity, compressed hierarchies, multisurfaces, and Heegaard splittings with
Cheeger regions. The final flowchart contains operations on these objects, not
on Khovanov complexes~\cite{lacktalk}.

The supplied July 2026 paper presents a hierarchy-based algorithm for
incompressibility and hence unknot recognition. Section~9 explicitly limits its
analysis to iteration counts in terms of maximal hierarchy length $L$ and
maximal pattern-complexity $g$. Proposition~9.1 bounds these iterations by
$L(g+1)^L$. That section does not bound both parameters in the input size and
states that the steps have not there been specified precisely enough for a
running-time estimate~\cite[Section~9]{lackpaper}. Thus the supplied talk's
announcement and the supplied paper's detailed algorithm should not be silently
identified with a fully specified implementation having the announced bound.

Here is the quantitative obstacle to transferring that bound to code. Even a
hypothetical $g=O(n^2)$ only makes
\[
 L(g+1)^L=\exp\!\bigl(O(\log L+L\log n)\bigr).
\]
This is quasi-polynomial if an applicable effective length is $O(\log n)$; with
$L=O(n)$ the same expression is $\exp(O(n\log n))$. Moreover, an iteration count
must be multiplied by the bit cost of constructing and transforming the
compressed states. Substituting a desired length bound without implementing
and justifying the acceleration mechanism would not establish a complexity
bound.

A genuine implementation of the advertised route would need compatible
representations and executable routines for handle structures with boundary
patterns, compressed normal surfaces and parallelity information, cutting and
simplifying these structures, hierarchical multisurfaces, and the required
Heegaard/Cheeger operations, with representation-size invariants connecting
these routines. The supplied \code{fast/} code has none of these data types.
Its frontier matchings are two-dimensional smoothing data, not compressed
three-manifold hierarchies. Renaming them ``multisurfaces'' would not bridge the
gap. Agol--Hass--Thurston's compression techniques provide relevant machinery
for such a different implementation~\cite{aht}, but are not a drop-in bound on
this scanner.

Our implementation therefore retains the complete existing decision route and
makes the improvements below. The restricted quasi-polynomial corollary in
Section~\ref{sec:bounds} follows from actual executable factorization, not from
an assumed hierarchy routine. This assessment concerns the provided materials
and this implementation; it is not a claim that the announced general result
is false.

\section{Execution architecture and exactness contract}

The default decision path is:
\begin{center}
\small
validate $\longrightarrow$ R1/R2 reduction $\longrightarrow$ descending test\\[3pt]
$\longrightarrow$ modular Alexander $\longrightarrow$ capped Jones evaluations\\[3pt]
$\longrightarrow$ exact Alexander $\longrightarrow$ factored Khovanov ranks.
\end{center}
Each successful simplification or positive unknot test is exact. Each invariant
filter may only return \KNOTTED. An inconclusive invariant result proceeds to
the next stage. The symbolic Alexander stage is retained both for compatibility
and to handle modular collisions. It is not required for completeness once the
Khovanov stage is present.

The Jones state cap defaults to 4096 and is a performance policy, not a
topological assumption. Exceeding it skips that filter. In contrast, exhausting
the overall time or homological-object budget prevents the requested complete
calculation and returns \UNKNOWN. The time budget is cooperative: it is polled
at stage boundaries, scan construction steps, determinant pivots, factorization
steps, and inside cancellation updates. A single expensive arithmetic operation
can still overrun the requested wall time. An object cap is not a byte-level
memory cap.

For the final criterion, Kronheimer--Mrowka prove that reduced Khovanov rank one
detects the unknot~\cite{km}. The implementation uses $\F_2$: by the universal
coefficient theorem, rational rank is no greater than $\F_2$ dimension. The
rational reduced homology has nonzero Euler characteristic specialization and
therefore positive dimension. Thus $\dim_{\F_2}\rKh(K)=1$ forces rational rank
one, and the detector applies. The converse follows from the explicit unknot
complex. Section~\ref{sec:splitting} explains why halving the unreduced rank is
valid over $\F_2$, including degree by degree.

\section{Exact one-sided modular rejection filters}

\subsection{Evaluating the Alexander minor without polynomial elimination}

Let $F_D(t)\in\Z[t]$ be the determinant of the same $(n-1)\times(n-1)$ Fox
minor used by the original implementation. Its entries have degree at most one,
so
\[
 \deg F_D\le n-1.
\]
For a knot this minor represents the Alexander polynomial up to multiplication
by a Laurent unit. In particular, if $D$ represents the unknot, then
\begin{equation}\label{eq:alexunit}
 F_D(t)=\varepsilon t^k,\qquad
 \varepsilon\in\{-1,1\},\quad 0\le k\le n-1.
\end{equation}
The polynomial matrix need never be constructed to test this necessary
condition at a fixed value.

\begin{proposition}[Sound modular rejection]\label{prop:alex}
Fix a prime $p$ and $a\in\F_p^\times$. If
\[
 F_D(a)\notin \{\pm a^k:0\le k\le n-1\}\subseteq\F_p,
\]
then $D$ is not the unknot.
\end{proposition}
\begin{proof}
Evaluation modulo $p$ is a ring homomorphism. If $D$ were the unknot,
\eqref{eq:alexunit} would put its evaluation in the displayed set. The
contrapositive proves the claim. No assumption about avoiding collisions is
needed for this implication.
\end{proof}

The implementation uses $p=1000000007$ and $a=-1,2$, in that order. At $-1$ the
allowed set is just $\{-1,1\}$; many knots are rejected at this first evaluation.
At $2$ the whole allowed set is retained, rather than incorrectly demanding the
value be $\pm1$. A monomial $t^k$ with $k>0$ is still a unit in the Laurent ring.

The numeric Fox matrix is built directly over $\F_p$. Modular Gaussian
elimination costs $O(n^3)$ field operations and $O(n^2)$ entries of storage.
Only nonzero entries of the current pivot row need updates. The prime is fixed,
so there is no coefficient swell. A zero determinant is also outside the unit
set and is a valid rejection. When $n=1$, the minor is empty and its determinant
is one, as required. A zero-crossing diagram is handled before this test.

\subsection{Jones evaluations as a matching-valued transfer computation}

A PD row $(a,b,c,d)$ has smoothings
\[
 0:(a,b)(c,d),\qquad 1:(a,d)(b,c).
\]
For $A\ne0$, put $\delta=-A^2-A^{-2}$. If $s\in\{0,1\}^n$ and $c(s)$ is the
number of circles in its complete smoothing, define the raw state sum
\begin{equation}\label{eq:raw}
 Z_D(A)=\sum_s A^{n-2|s|}\delta^{c(s)}.
\end{equation}
This convention assigns $\delta$ to a crossing-free circle. With the compatible
writhe convention, the normalized bracket is
\begin{equation}\label{eq:normalized}
 J_D(A)=\delta^{-1}(-A^3)^{-w(D)}Z_D(A).
\end{equation}
It equals one on the unknot. These are the usual bracket-state construction and
writhe correction, with an explicit normalization chosen to match the input
ports~\cite{kauffman}.

The important acceleration is aggregation before homology. At any partial scan,
store one coefficient vector per boundary matching, rather than a chain complex
with separate objects in all degrees. Attaching one smoothing closes $c$ circles
and changes the boundary matching. Multiply the coefficient by $A$ or $A^{-1}$
and by $\delta^c$, and sum all contributions having the same resulting matching.
The implementation evaluates at $A=2$ and $A=3$ simultaneously in $\F_p$.

\begin{proposition}[Frontier recurrence]
After any prefix of the scan, the coefficient of a boundary matching is the sum
of the partial-resolution weights that produce that matching, including the
weights of all circles already closed. At the empty final boundary the
coefficient is $Z_D(A)$.
\end{proposition}
\begin{proof}
Initially there is one empty partial resolution of weight one. A new crossing
partitions extensions into its two smoothings. Each extension multiplies the
old weight by the appropriate $A^{\pm1}$ and by one $\delta$ for each new closed
circle. Future gluing depends only on the remaining matching, so summing
coefficients for identical matchings loses no information. Induction gives the
invariant. At the end every smoothing circle has closed and all crossing weights
have been included, giving \eqref{eq:raw}.
\end{proof}

Both chosen evaluation points have nonzero $\delta$. Therefore all inverses in
\eqref{eq:normalized} exist. If either normalized value differs from one, the
knot is nontrivial. If both equal one, the result is inconclusive. In particular,
the program does not assume that the Jones polynomial detects the unknot, nor
that two evaluations determine the polynomial.

Let $S$ be the largest retained state count. Apart from ordering, the transfer
algorithm costs $\poly(n,w)S$ for maximum boundary size $w$, with two evaluations
a constant factor. For an arbitrary scan, a safe combinatorial bound on the
number of matchings is $(w-1)!!=2^{O(w\log w)}$, with the empty matching counted
as one. A Catalan bound requires an appropriate single-disc cyclic-boundary
condition; the greedy implementation does not establish that condition. It
would be unsound to substitute it merely because the original link diagram is
planar. With the default fixed state cap, the filter has only polynomial
attempted work before it either completes or falls through.

\subsection{Writhe and independently checkable rejection records}

The original \code{signs()} routine gives local Fox-calculus signs relative to
the under-port direction $a\to c$. These are not in general the oriented
crossing signs. The new \code{writhe()} follows the actual traversal and uses the
incoming under and over ports together. In its convention, pairs $(0,3)$ and
$(2,1)$ have sign $+1$; the other two possibilities have sign $-1$. A positive
braid generator has writhe $-1$ in this particular input convention. Rotating a
PD row by two ports must leave the writhe and normalized bracket unchanged; this
is a regression test. Replacing the Fox signs themselves would instead have
changed the existing Alexander matrix convention, so they are retained.

A modular rejection result contains the reduced PD, the R1/R2 reduction trace,
and the exact field witness. \code{verify_rejection} first replays the legal
moves on the original input, checks the resulting diagram, then recomputes the
minor or bracket and checks the inequality. Tampering with the witness, prime,
input diagram, or reduction trace is rejected. This is a reproducible certificate
checker, not a formally verified kernel and not an independent implementation
of the invariant algorithm. It intentionally does not claim to verify arbitrary
Khovanov output transcripts.

\section{Faster cobordism arithmetic and cancellation}

\subsection{The retained local algebra}

The scanner is based on Bar-Natan's local cancellation approach~\cite{bn}.
Boundary objects are perfect matchings. A morphism is an $\F_2$ combination of
dot-subsets of the circles obtained by gluing the source and reversed target
matchings. Delooping a closed circle replaces it with two summands. Arithmetic
uses the Frobenius algebra
\[
 A_0=\F_2[x]/(x^2),\qquad
 \Delta(1)=1\otimes x+x\otimes1,\quad \Delta(x)=x\otimes x,
 \quad \epsilon(1)=0,\quad\epsilon(x)=1.
\]
The geometric composition routine in the baseline is retained as a reference
path. Its repeated reconstruction of glued surfaces is avoidable when the same
local data recur many times in the complex.

\subsection{Every unit is its own inverse}

For a matching $m$ with $r$ arcs, its endomorphism ring is
\[
 R_m\simeq\F_2[x_1,\ldots,x_r]/(x_1^2,\ldots,x_r^2).
\]
The baseline already recognizes units by a nonzero constant term. The new code
uses the following stronger simplification.

\begin{lemma}\label{lem:inverse}
If $f\in R_m$ has constant coefficient one, then $f^{-1}=f$.
\end{lemma}
\begin{proof}
Write $f=1+u$, where $u$ has zero constant coefficient. Every positive-degree
square-free monomial squares to zero. In characteristic two, cross terms in the
square of their sum cancel, so $u^2=0$. Thus $f^2=(1+u)^2=1$.
\end{proof}

This argument applies to arbitrary combinations of monomials; it does not depend
on retaining a quantum grading or on $f$ being homogeneous. In the same ring,
multiplication is particularly cheap: two dot-subsets with an intersection
multiply to zero; otherwise they multiply to their union, with duplicates
cancelled modulo two. The implementation uses this path when all three
matchings in a composition agree. Zero and identity morphisms have immediate
shortcuts.

\subsection{Bounded memoization and immutable coefficients}

Morphism values used as memoization keys and cached results are frozen. General
composition is cached by both morphisms and all three boundary matchings. The
entire local crossing transformation is cached by the input morphism, source and
target matchings, smoothing choices, boundary, and crossing ports. This avoids
recomputing the same geometric/delooping transformation for different copies of
a homological object. A fresh outer dictionary is returned where the calling
code expects mutable indexing, while cached coefficient values remain immutable.

The entry limits are 16384 for circle decompositions, 16384 for smoothing gluing,
65536 for composition, and 32768 for crossing transformations. These are entry
bounds, not byte bounds: a large morphism is still large. Internal cached geometry
is treated as read-only. \code{clear_caches()} and \code{cache_info()} support
controlled experiments and release of retained state between runs.

\subsection{Sparse, lazy Markowitz-style cancellation}

Suppose an invertible differential entry $\phi:b\to c$ is selected. For each
remaining incoming $\delta:a\to c$ and outgoing $\gamma:b\to f$, cancellation
replaces the direct entry by
\begin{equation}\label{eq:schur}
 d'_{af}=d_{af}+\gamma\phi^{-1}\delta.
\end{equation}
The minus sign of the general Schur complement is a plus over $\F_2$. The pair
$b,c$ is then deleted. Bar-Natan's cancellation lemma identifies the resulting
complex as chain-homotopy equivalent to the original~\cite[Lemma~4.2]{bn}.

The number of potential updates associated with a pivot is
\[
 \mu(b,c)=\bigl(|\operatorname{in}(c)|-1\bigr)
          \bigl(|\operatorname{out}(b)|-1\bigr).
\]
The baseline uses a LIFO list of available units. The optimized default instead
uses a heap with priority $(\mu,\text{degree sum},b,c)$. On extraction it checks
that both objects and the unit still exist, recomputes the current priority,
and reinserts if stale. New invertible entries created by
\eqref{eq:schur} are inserted.

This is deliberately described as a \emph{lazy Markowitz-style heuristic}, not
an exact global minimum-degree or minimum-fill algorithm. Some entries whose
degrees have decreased can retain stale larger keys until popped. That affects
performance, not validity: every actually cancelled pivot is an extant unit and
therefore supports the same exact cancellation lemma. Between cancellations,
only finitely many stale entries require refreshing. A cancellation removes two
objects, so the process terminates. A persistent unit cannot be lost: it was
present initially or inserted when created, and any stale live copy is refreshed
rather than silently discarded.

Caching and pivot selection solve different problems. Caching makes a repeated
update cheap; pivot selection can avoid millions of updates altogether. The
large five-braid ablation in Section~\ref{sec:experiments} separates these two
effects. The original \code{compositions} statistic is kept for comparability:
it increments twice per incoming/outgoing update pair. It is not a count of
expensive cache misses or necessarily of actual geometric reconstructions.

\section{Certified diagrammatic connected-sum factorization}
\label{sec:factors}

\subsection{Recognizing a valid cut without a prime-decomposition oracle}

Traverse the diagram once and record the double-occurrence word of its crossing
indices. Search proper cyclic intervals in which every crossing appearing in
the interval appears twice. Such an interval has exactly two traversed cut
edges. The implementation also computes the faces of the planar rotation system
and requires the cut edges to have the same two \emph{distinct} incident faces.
It favors a balanced crossing partition among the cuts found.

\begin{proposition}[Sound diagrammatic splitting]
A cut satisfying these tests yields a connected-sum decomposition of the knot.
Identifying the two cut-edge labels separately in each side gives valid PDs for
the factors, with the original crossings partitioned between them.
\end{proposition}
\begin{proof}
The interval condition assigns every crossing wholly to one side and gives two
edges joining the two parts. The face condition makes their dual edges a simple
two-edge cycle: one arc can be drawn in each of the two distinct faces, joining
their intersection points with the cut edges. Their union is a Jordan curve on
the projection sphere, avoids all vertices, and meets the diagram transversely
in exactly those two points. It separates the two crossing sets. Pushing this
curve to a sphere in three-space gives the usual connected-sum sphere meeting
the knot twice. Closing each resulting arc across its disc corresponds exactly
to identifying the two cut labels in its factor PD. The code additionally
validates both resulting PDs.
\end{proof}

The procedure is conservative. It does not find every topological connected-sum
sphere after arbitrary diagram changes and is not a prime-decomposition
algorithm. Failure to find a cut means only that this implementation will scan
that piece directly. An explicit stack processes the decomposition tree without
Python recursion-depth dependence. Each cut records the local PD, the two edge
labels, the local crossing partitions, and their relation to the original
crossing indices.

For a node with $m$ crossings, each of $2m$ interval starts can be scanned in
$O(m)$ set operations, giving $O(m^2)$ word-RAM work. There are fewer than $2n$
nodes, so $O(n^3)$ is a conservative total bound. Favoring balance helps in
practice but is not used to assert that every decomposition tree is balanced.

\subsection{Why the rank product and degree convolution are exact}
\label{sec:splitting}

We include the algebraic details because it would be incorrect to multiply
unreduced ranks directly, or to halve arbitrary coefficient-field ranks.
Let $C(D)$ be the cube complex over $\F_2$, with the basepoint away from
crossings, and let $X$ multiply by $x$ on the circle through that basepoint.
The reduced complex is $\widetilde C(D)=\operatorname{im}X=\ker X$.

Let $\nu$ act on each enhanced state by the sum of the operations replacing one
$x$ label by $1$. A local check against multiplication and comultiplication
gives
\[
 \nu d=d\nu,\qquad Xd=dX,\qquad \nu X+X\nu=1.
\]
For example, applying the total derivation to
$\Delta(1)=1\otimes x+x\otimes1$ gives two copies of $1\otimes1$, hence zero;
on $\Delta(x)$ it gives $\Delta(1)$. These identities are the standard
characteristic-two splitting mechanism~\cite[Section~3.2]{shumakovitch}.

\begin{lemma}[Unreduced/reduced splitting]
After forgetting quantum grading, $C(D)$ is the direct sum of two copies of
$\widetilde C(D)$ as a complex preserving cube degree. Consequently
\[
 \dim\Kh^h(D;\F_2)=2\dim\rKh^h(D;\F_2).
\]
\end{lemma}
\begin{proof}
For $r\in\ker X$, the identity $X\nu r=r$ shows that
$\nu:\widetilde C\to C$ is a chain-map section of $X:C\to\widetilde C$.
Every $c\in C$ decomposes as
$c=(c+\nu Xc)+\nu Xc$, with the first term in $\ker X$.
The two summands intersect trivially since $X\nu$ is the identity on
$\widetilde C$. Both maps preserve cube degree. Taking homology proves the rank
identity. Quantum shifts distinguish the two copies, but are not being reported.
\end{proof}

\begin{lemma}[Connected-sum product]
For diagrammatic factors $D_1,D_2$ with their basepoints at the joining arcs,
\[
 \widetilde C(D_1\#D_2)\cong
 \widetilde C(D_1)\otimes_{\F_2}\widetilde C(D_2)
\]
in raw cube degree, after forgetting quantum grading. Hence their reduced rank
polynomials multiply.
\end{lemma}
\begin{proof}
A smoothing of the connected sum is a pair of smoothings with the two marked
circles joined. In the reduced complex the joined circle is labelled $x$.
This gives a basis bijection with pairs of enhanced states each having marked
label $x$. At every crossing the saddle occurs in one factor only. The
multiplication and comultiplication formulas above show that this bijection
intertwines the differentials, including saddles involving the marked circle.
The tensor differential has no sign distinction over $\F_2$. The K\"unneth
formula over a field now identifies homology with the tensor product of factor
homologies. Since the crossing sets are partitioned, cube degrees add.
\end{proof}

Write $H_i(z)=\sum_h\dim\Kh^h(D_i;\F_2)z^h$. The implementation computes
\begin{equation}\label{eq:convolution}
 H_D(z)=2\prod_i \frac{H_i(z)}{2}.
\end{equation}
The division is coefficientwise and checked for evenness. Multiplication is
ordinary sparse integer convolution. Identical normalized factor PDs share a
memoized rank result. The program never builds the tensor-product generators.
It returns an exact integer total and a degree histogram with at most $n+1$
entries.

\subsection{Bounds and a genuine asymptotic improvement}
\label{sec:bounds}

\begin{theorem}[Complexity of the delivered algorithm]
Let $n_i$ be the crossing counts of the diagrammatic factors actually found, and
$b=\max_i n_i$. The uncapped exact rank routine has a conservative bound
\[
 T(D)\le\poly(n)+\sum_i 2^{O(n_i)}
        \le\poly(n)+n\,2^{O(b)}.
\]
The default complete recognizer has the same form after its polynomial
preprocessing and fixed-cap filters. In particular, the general bound is
$2^{O(n)}$.
\end{theorem}
\begin{proof}
Consider an unfactored piece of $m$ crossings. Its explicit resolution cube has
$2^m$ states and at most $m+1$ smoothing circles per state, hence at most
$2^{2m+1}$ enhanced states. Delooping and scanning cannot increase the number of
objects beyond an exponential-size expanded state description; cancellation
only deletes objects. There are exponentially many possible dotted basis terms
at a boundary of $O(m)$ points, so every morphism can be stored and composed in
$2^{O(m)}$ work. The number of cancellation updates is bounded by a polynomial
in the exponential number of objects and entries. Lazy heap operations and
bounded memoization do not change this exponential upper bound.

The decomposition is polynomial. Factor calculations are sequential and sum to
$\sum_i2^{O(n_i)}$. Intermediate histogram coefficients have $O(n)$ bits because
all dimensions are bounded by an exponential-size cube. There are $O(n)$ degrees
and $O(n)$ factors; even straightforward integer convolution therefore costs a
polynomial number of bit operations. This proves the first bound, and at most
$n$ nonempty factors give the second. Modular elimination is polynomial; the
fixed-cap Jones attempt is polynomial; the retained symbolic Alexander
calculation has polynomial degree and coefficient-bit growth for its linear
integer matrix entries. These stages do not increase the general exponential
bound.
\end{proof}

\begin{corollary}
On diagrams for which the found maximum factor size is $O(\log n)$, the delivered
algorithm runs in polynomial time. If that maximum is $O(\log^2 n)$, it runs in
$n^{O(\log n)}$ time.
\end{corollary}

This is a bound on \emph{found diagrammatic factors}, not merely on the abstract
prime factors of the knot. It does not provide a quasi-polynomial bound for a
large prime knot or a diagram hiding its connected-sum structure.

The family
\[
 D_k=\widehat{\sigma_1^3\sigma_2^3\cdots\sigma_k^3}
 \quad\text{on }k+1\text{ strands}
\]
is a concrete exponential-to-polynomial improvement for the \emph{full rank}
API. It is a visible connected sum of $k$ trefoils, has $n=3k$, and has reduced
rank $3^k$. The baseline ends with $2\cdot3^k$ separate homology objects and
therefore takes at least $\Omega(3^k)$ object work. The new routine finds
three-crossing factors and evaluates \eqref{eq:convolution} compactly in
polynomial time. This does not claim an exponential speedup of default
\emph{recognition} on that family: the baseline can already reject these knots
by Alexander polynomial.

The same family shows why scan boundary size alone is not a valid complexity
parameter for this implementation. It admits the measured width-four scan, yet
has exponentially growing homological multiplicity. Boundary matching types and
the number of copies of each type are different quantities. The original
width-only explanation is therefore explicitly not repeated as a guarantee.

\section{Faster polynomial preprocessing}

\subsection{An equivalent greedy order in \texorpdfstring{$O(n\log n)$}{O(n log n)} per start}

The baseline greedy rule repeatedly scans every unprocessed crossing, preferring
more shared boundary edges, then more self-loops, then the lower crossing index.
For one start this takes $O(n^2)$ comparisons. The new implementation builds the
crossing adjacency lists once. Each crossing has at most four incident edges.
When a crossing is processed, only its still-unprocessed neighbors have their
shared-edge scores increased. A lazy heap stores
\[
 (-\text{shared edges},-\text{self-loops},\text{index}).
\]
There are $O(n)$ score changes and heap entries, so a start costs
$O(n\log n)$ time and $O(n)$ space. The original start sequence and tie-breaking
are retained exactly, including the baseline's behavior when the requested
number of starts does not divide $n$. Randomized differential tests compare
\emph{every} start against the unchanged baseline. Thus improvements in the raw
scan comparisons are not explained by silently choosing easier crossing orders.

\subsection{All descending basepoints by one interval sweep}

Fix an oriented traversal of its $2n$ crossing encounters. A diagram is descending
from a basepoint if the overpass is the first encounter with every crossing.
For a crossing with over-encounter $o$ and under-encounter $u$, the bad starting
positions form the cyclic interval $(o,u]$: those starts see $u$ first. Add one
to each such interval using a cyclic difference array. A prefix sum gives the
number of bad crossings at every start. Exactly the zero positions are valid
descending basepoints.

Repeat for the reversed traversal and choose the smallest valid dart, matching
the previous API. This takes $O(n)$ work and space instead of retraversing the
entire diagram from every candidate dart. The fact that a descending knot
diagram represents the unknot is unchanged. The R1/R2 reduction strategy itself
is not redesigned; deadline checks are added between its decreasing moves.

\section{Validation and its limits}

The delivered regression suite has 30 test methods. The original 19 methods are
retained; only expectations of default rejection-method names and explicit
filter-disabling options are adapted where the pipeline has changed. The exact
unchanged original tests remain in \code{baseline/tests/} and also passed.
The new tests cover algebraic identities, randomized oracle comparisons,
scan-order equivalence, descending-basepoint equivalence, factor products and
degree histograms, port rotations, label permutations, invalid orders, resource
budgets, cache bounds, and tampered certificates. The original difficult
Alexander-trivial named examples are retained.

A separate exhaustive script enumerates every word of length 2, 4, or 6 over
$\{\sigma_1^{\pm1},\sigma_2^{\pm1}\}$ whose closed three-braid has one component.
These are diagram words, not distinct knot isotopy classes. For each word it
compares the direct scanner and factored scanner with an independent dense
reduced cube, checks $d^2=0$ during the direct scan, compares the Jones transfer
calculation to a full independent smoothing-state sum, and checks the default
recognizer's verdict. For every modular rejection it recomputes the certificate.

\begin{center}
\begin{tabular}{rr}
\toprule
Word length & One-component words checked\\
\midrule
2 & 8\\
4 & 160\\
6 & 2688\\
\midrule
Total & 2856\\
\bottomrule
\end{tabular}
\end{center}
There are 1712 unknot words and 1144 rejected nontrivial words in this set, with
zero discrepancies. The dense cube reference is adapted from the original
archive's independent test oracle. It enumerates enhanced states, explicitly
builds saddle differentials, and performs $\F_2$ column elimination without
using the scanner's cobordism composition or gluing. The Jones oracle separately
uses an edge-disjoint-set computation for each complete smoothing, without the
frontier \code{glue} function.

The large 36-crossing example is beyond the dense reference run performed here.
Its two completed cancellation strategies agree on the \emph{entire} raw
homological degree histogram. Six checked R2 reductions give a 24-crossing
diagram whose optimized calculation, with $d^2$ checks, again has reduced rank
2949. Its independently computed exact Alexander determinant is also 2949.
These are useful consistency checks; the determinant agreement is not an
independent full-rank proof and is not reported as one. No Lean formalization
or third-party full-rank computation is claimed.

\section{Same-machine experiments}
\label{sec:experiments}

\subsection{Methodology}

The measurements used CPython 3.13.5, Linux x86-64, and an AMD EPYC 9V74 processor.
Each sample runs in a fresh subprocess, pinned to the same eligible logical CPU
when the operating system supports affinity, with \code{PYTHONHASHSEED=0}.
Samples are serial. Module import, input JSON parsing, and initial conversion
into a diagram are outside the timed region. The public API's own validation,
ordering, and any enabled factorization are inside it. Caches are cleared before
timing. Engine order is alternated between repetitions. Except for the long
five-braid case, tables show medians of three samples. Millisecond-scale ratios
should not be interpreted as high-precision measurements of hardware-independent
constants.

The baseline is the actual supplied source, not a reconstructed approximation.
The \code{rank} task uses no R1/R2 simplification. Factoring is disabled in generic
raw-rank comparisons and enabled only in the connected-sum suite. Completed rank
runs must agree degree by degree; completed recognition runs must agree in
verdict. Timeouts are censored observations, not fabricated finishing times.
The harness records both the cooperative budget and any parent-process watchdog
expiration. Source fingerprints are recorded at measurement time; the final
manifest separately identifies the delivered files.

\subsection{Raw homology calculations}

\begin{center}
\small
\setlength{\tabcolsep}{5pt}
\begin{tabular}{lrrrrr}
\toprule
Instance & $n$ & Reduced rank & Baseline (ms) & New (ms) & Speedup\\
\midrule
Conway & 11 & 33 & 34.402 & 11.951 & 2.88$\times$\\
Kinoshita--Terasaka & 11 & 33 & 38.206 & 13.160 & 2.90$\times$\\
Hard unknot & 8 & 1 & 2.321 & 2.754 & 0.84$\times$\\
$T(3,5)$ & 10 & 7 & 9.731 & 9.599 & 1.01$\times$\\
Unknot braid & 40 & 1 & 8.290 & 3.494 & 2.37$\times$\\
$T(2,41)$ & 41 & 41 & 82.195 & 25.089 & 3.28$\times$\\
Random four-braid & 21 & 25 & 259.337 & 81.672 & 3.18$\times$\\
Random four-braid & 31 & 45 & 301.728 & 70.107 & 4.30$\times$\\
Random four-braid & 41 & 109 & 954.437 & 249.370 & 3.83$\times$\\
Random six-braid & 31 & 85 & 515.203 & 88.843 & 5.80$\times$\\
\bottomrule
\end{tabular}
\end{center}

Most nontrivial examples improve by roughly three to six times, while the tiny
hard-unknot raw-rank case regresses. This is consistent with saving repeated
algebra and updates where there is enough work to amortize heap, validation,
and caching overhead. No universal constant-factor speedup is claimed.

\subsection{The large five-braid: separating caching from cancellation}

The exact word, also saved as \code{examples/random_5_braid_36.json}, is
\begin{lstlisting}
[-4,-3,4,4,-3,2,1,1,-1,-1,3,-4,-1,3,2,2,-2,-1,
 -3,2,3,-3,-4,-4,-3,3,1,3,-2,-4,-4,3,-3,3,-4,-2]
\end{lstlisting}
All rows below scan this same original 36-crossing diagram, with the same
crossing order, without factoring or Reidemeister simplification.

\begin{center}
\small
\begin{tabular}{lrrr}
\toprule
Engine & Seconds & Reduced rank & Update counter\\
\midrule
Supplied baseline & $>60$ & not obtained & not completed\\
New arithmetic/caches, LIFO & 49.576 & 2949 & 71\,475\,304\\
New arithmetic/caches, lazy fill-aware & 1.110 & 2949 & 190\,930\\
\bottomrule
\end{tabular}
\end{center}
The baseline raised its time-budget exception after 60.073 seconds; it did not
complete in the budget. The resulting conservative comparison is more than
$60/1.110>54$ times faster within this test budget. The completed LIFO ablation
versus the new pivot order gives about $44.7$ times improvement, with all other
optimized machinery held fixed. Each of these expensive configurations was
measured once.

Both completed scans have the same maximum pre-cancellation object count
17694, maximum post-cancellation count 5898, 16884 eliminations, and boundary
width eight. What changes dramatically is the number of attempted cross updates,
by a factor of about 374.4. The entire rank histogram agrees. Process memory
high-water marks, including interpreter/import overhead, are approximately
276.2 MiB for cached LIFO and 127.1 MiB for the new pivot heuristic. The baseline's
censored peak is not its eventual peak and is not used as a completed-run memory
comparison.

\subsection{Default recognition, including regressions}

\begin{center}
\small
\setlength{\tabcolsep}{5pt}
\begin{tabular}{lrrrr}
\toprule
Instance & $n$ & Baseline (ms) & New (ms) & Speedup\\
\midrule
Conway & 11 & 34.036 & 1.610 & 21.14$\times$\\
Kinoshita--Terasaka & 11 & 38.233 & 1.563 & 24.47$\times$\\
Hard unknot & 8 & 1.915 & 2.735 & 0.70$\times$\\
$T(3,5)$ & 10 & 0.661 & 0.266 & 2.48$\times$\\
Trefoil & 3 & 0.097 & 0.119 & 0.82$\times$\\
Unknot braid & 40 & 6.889 & 6.995 & 0.98$\times$\\
Determinant-one grid & 10 & 0.808 & 0.281 & 2.87$\times$\\
Scrambled unknot grid & 6 & 0.302 & 0.283 & 1.07$\times$\\
\bottomrule
\end{tabular}
\end{center}

Conway and Kinoshita--Terasaka have trivial Alexander polynomials in the supplied
examples. The new exact Jones evaluations reject them without constructing a
homological differential, explaining the roughly 21--24-fold gains. The hard
unknot reaches the complete fallback and pays for filters that do not reject it;
it is about 43 percent slower on this small input. The tiny trefoil also shows a
small overhead. The table includes these regressions rather than only the wins.

\subsection{Compact products instead of exponentially many objects}

\begin{center}
\small
\setlength{\tabcolsep}{5pt}
\begin{tabular}{rrrrrr}
\toprule
Factors & Crossings & Baseline (ms) & New (ms) & Peak objects & Reduced rank\\
\midrule
2 & 6 & 2.562 & 1.911 & 18 & $3^{2}$\\
4 & 12 & 15.440 & 2.611 & 18 & $3^{4}$\\
6 & 18 & 176.910 & 3.008 & 18 & $3^{6}$\\
8 & 24 & 1305.710 & 3.638 & 18 & $3^{8}$\\
20 & 60 & not run & 6.618 & 18 & $3^{20}$\\
50 & 150 & not run & 22.922 & 18 & $3^{50}$\\
100 & 300 & not run & 65.597 & 18 & $3^{100}$\\
\bottomrule
\end{tabular}
\end{center}

For eight trefoil factors, the old rank routine peaks at 39366 pre-cancellation
objects and ends with 13122 unreduced homology objects. The new routine needs
at most 18 pre-cancellation objects in a factor. At 100 factors it still uses
that same object peak and returns
\[
 3^{100}=515377520732011331036461129765621272702107522001.
\]
The input has 300 crossings. A full degree histogram, not just the scalar
product, is computed. Baseline runs above eight factors were not performed in
this experiment and are marked ``not run,'' not timed out or assigned an
extrapolated time.

\subsection{Ordering alone}

\begin{center}
\small
\setlength{\tabcolsep}{5pt}
\begin{tabular}{rrrr}
\toprule
$n$ & Baseline (ms) & New (ms) & Speedup\\
\midrule
128 & 3.995 & 0.199 & 20.1$\times$\\
256 & 15.995 & 0.365 & 43.8$\times$\\
512 & 65.135 & 0.725 & 89.9$\times$\\
1024 & 251.053 & 1.629 & 154.1$\times$\\
\bottomrule
\end{tabular}
\end{center}

These inputs are closures of $\sigma_1\cdots\sigma_n$ on $n+1$ strands. The timed
task is a single greedy order, not knot recognition. The output orders agree
exactly. The observed scaling is consistent with the proved replacement of
quadratic rescanning by incremental heap updates, but the complexity claim rests
on the operation count, not on fitting these four timings.

\section{Using and reproducing the deliverable}

No installation is required. From the archive root:
\begin{lstlisting}
cd fast
python -m fastunknot recognize examples/conway.json --output conway-proof.json
python -m fastunknot verify examples/conway.json conway-proof.json
python -m fastunknot khovanov examples/random_5_braid_36.json --no-factors
python -m fastunknot recognize examples/hard_unknot_8.json
\end{lstlisting}
Python 3.10 or later is required. For a normal installed console command,
\code{python -m pip install ./fast} can be run from the archive root. There are no
runtime package dependencies. The same source commands work in PowerShell.

For a controlled fallback comparison, disable both invariant families:
\begin{lstlisting}
python -m fastunknot recognize examples/conway.json --no-alexander --no-jones
\end{lstlisting}
The \code{--no-alexander} option disables modular and symbolic Alexander, but not
Jones. The \code{khovanov} subcommand bypasses all recognition preprocessing;
\code{--pivot lifo} selects the cached old pivot strategy. An explicitly supplied
library scan order must be a full permutation and disables factorization.

From the archive root, the validation and measurement commands are:
\begin{lstlisting}
python -m unittest discover -s fast/tests -v
python tools/validate_exhaustive.py
python tools/benchmark.py
\end{lstlisting}
The last command takes several minutes because it deliberately repeats the
baseline time-limited case and cached-LIFO ablation. \code{--only conway} or
\code{--only ordering_} limits the requested cases. The exhaustive script writes
the case counts, rank histogram, certificate count, result-stream digest, and
large-input consistency checks to \code{results/exhaustive.json}.

For library use:
\begin{lstlisting}
from fastunknot import Diagram, recognize, khovanov_rank

d = Diagram.from_braid(3, [1, -2, 1, -2])
answer = recognize(d, seconds=30, max_objects=200_000)
print(answer.status, answer.method)

homology = khovanov_rank(d.pd, factor_connected=True)
assert homology["reduced_rank"] == 5
\end{lstlisting}
The public JSON contracts retain the distinction between \UNKNOT, \KNOTTED,
and \UNKNOWN. A completed knot verdict has exit code zero; a resource limit has
code three; invalid input or options have code two. Certificate verification
uses zero for a verified modular rejection and one for failure.

The report can be rebuilt from \code{report/} with two runs of
\code{pdflatex -interaction=nonstopmode -halt-on-error unknot_speedup.tex}.
No bibliography download, shell escape, external figure, or network connection
is needed. \code{changes.patch} exposes the implementation changes relative to
the original source. The original 600-second Windows benchmark is retained only
under \code{baseline/results/}; all new comparisons are under \code{results/}.

\section{Remaining limits and concrete next implementation targets}

The generic worst case is still exponential. Diagrammatic factors can all be
large; the Jones filter can be inconclusive or exceed its cap; Alexander can be
trivial for a nontrivial knot; the homological complex can have exponentially
many surviving objects; and bounded caches do not bound total RAM. The lazy
pivot heuristic is not optimal and can lose on particular inputs. The test suite
is extensive but finite, and the certificate checker shares its invariant
implementation with the producer.

Further practical work can target a compact integer encoding of matchings and
dot monomials, moving cache-key construction and finite-field elimination into
compiled kernels, and improved separator-based order heuristics. Those changes
would need new ablations because changing crossing order and cancellation order
simultaneously can obscure the source of a gain. A complete general
quasi-polynomial hierarchy implementation is a distinct target: it requires the
three-manifold representations and compressed-operation proofs identified in
Section~2, not just a faster scalar matrix reducer.

The present deliverable is therefore a tested, usable speed-up with exact
one-sided filters, a much faster difficult raw-rank example, and a genuine
exponential-to-polynomial improvement for compactly decomposable rank problems.
Its bounds and evidence are deliberately stated at the level actually supported
by the source code and measurements.

\begin{thebibliography}{9}
\bibitem{lacktalk}
Marc Lackenby, \emph{Unknot recognition in quasi-polynomial time}, talk slides,
2021. Supplied as \code{docs/quasipolynomial-talk.pdf}; 109 physical pages.
The announcement is on physical pages 13--14; speed-up ingredients on page 72;
final flowchart on page 109.
\url{https://people.maths.ox.ac.uk/lackenby/quasipolynomial-talk.pdf}.

\bibitem{lackpaper}
Marc Lackenby, \emph{Incompressible surfaces, hierarchies and unknot recognition},
arXiv:2607.23350v1, 25 July 2026. Supplied TeX source and figures in
\code{docs/arXiv-2607.23350v1/}; the complete PDF is available from arXiv.
In particular, Section 9 and Proposition 9.1,
physical pages 34--35.
\url{https://arxiv.org/abs/2607.23350}.

\bibitem{bn}
Dror Bar-Natan, \emph{Fast Khovanov Homology Computations},
Journal of Knot Theory and Its Ramifications \textbf{16} (2007), no.~3,
243--255; arXiv:math/0606318.
\url{https://arxiv.org/abs/math/0606318}.

\bibitem{km}
P.~B. Kronheimer and T.~S. Mrowka, \emph{Khovanov homology is an unknot-detector},
arXiv:1005.4346 (2010).
\url{https://arxiv.org/abs/1005.4346}.

\bibitem{shumakovitch}
Alexander N. Shumakovitch, \emph{Torsion of the Khovanov homology},
Fundamenta Mathematicae \textbf{225} (2014), 343--364;
arXiv:math/0405474v2 (2018). In particular, the $X,\nu$ identities and
Corollary 3.2.C.
\url{https://arxiv.org/abs/math/0405474}.

\bibitem{kauffman}
Louis H. Kauffman, \emph{Combinatorial Knot Theory and the Jones Polynomial},
arXiv:2204.12104 (2022), the originator's exposition of the bracket state model.
\url{https://arxiv.org/abs/2204.12104}.

\bibitem{aht}
Ian Agol, Joel Hass, and William P. Thurston,
\emph{The Computational Complexity of Knot Genus and Spanning Area},
arXiv:math/0205057 (2002).
\url{https://arxiv.org/abs/math/0205057}.
\end{thebibliography}
\end{document}
